\documentclass[a4paper,12pt]{article}
\usepackage[utf8]{inputenc}
\usepackage[T1]{fontenc}
\usepackage[italian]{babel}
\usepackage{pgfplots}
\usepackage{amsmath}
\usepackage{float}
\usepackage{booktabs}
\usepackage{multirow}
\usepackage{graphicx}
\usepackage{geometry}
\geometry{margin=2cm}

\pgfplotsset{compat=1.18}

\title{Report 4.11 -- Risultati per $k$ con $TIMES = 10$}

\begin{document}

\maketitle

\newpage

\section{Grafici dei risultati per $k$ e $TIMES = 10$}

I grafici mostrano rispettivamente l'andamento di:
\begin{itemize}
    \item Avg fQ
    \item fQ Min Found
    \item Profit Found
    \item Weight Found
    \item Profit GAP
    \item Weight GAP
    \item n of fQ Min Found
\end{itemize}
al variare di $k \in \{1000, 2000, 3000, 4000, 5000\}$, confrontando
i risultati usando QALS e NON utilizzando QALS.

I risultati riportati sono stati ottenuti utilizzando:

\begin{itemize}
    \item $\lambda = \frac{\sum_{i=1}^{n} p_i}{3}$
    \item l'istanza \texttt{ksp\_3.txt}
    \item \texttt{TIMES = 10}
\end{itemize}

\begin{table}[H]
\centering
\begin{tabular}{lc}
\toprule
Profit ottimo & 2679 \\
Weight ottimo & 496 \\
TIMES & 10 \\
\bottomrule
\end{tabular}
\caption{Soluzione di riferimento del problema di knapsack.}
\end{table}

\section{Tabella riassuntiva}

\begin{table}[H]
\centering
\resizebox{\textwidth}{!}{%
\begin{tabular}{c|rrrrrrr|rrrrrrr}
\toprule
\multirow{2}{*}{$k$} & \multicolumn{7}{c|}{QALS} & \multicolumn{7}{c}{NO QALS} \\
\cmidrule(lr){2-8} \cmidrule(lr){9-15}
 & Avg fQ & fQ Min & Profit & Weight & Profit GAP & Weight GAP & n fQ Min
 & Avg fQ & fQ Min & Profit & Weight & Profit GAP & Weight GAP & n fQ Min \\
\midrule
1000 & 17350642347.4 & 8241260668.67 & 3423 & 2275 & -744 & -1779 & 1 & -714267512.7 & -714267662 & 2483 & 501 & 196 & -5 & 6 \\
2000 & 18615897688.77 & 12244772130.67 & 3481 & 2635 & -802 & -2139 & 1 & -714267601.6 & -714267662 & 2483 & 501 & 196 & -5 & 8 \\
3000 & 17038731132.6 & 12039123606.67 & 4212 & 2618 & -1533 & -2122 & 1 & -714267631.8 & -714267662 & 2483 & 501 & 196 & -5 & 9 \\
4000 & 20763157976.07 & 16020029113.67 & 4249 & 2926 & -1570 & -2430 & 1 & -714267662 & -714267662 & 2483 & 501 & 196 & -5 & 10 \\
5000 & 17864058550.1 & 6034515014.67 & 2873 & 2041 & -194 & -1545 & 1 & -714267662 & -714267662 & 2483 & 501 & 196 & -5 & 10 \\
\bottomrule
\end{tabular}%
}
\caption{Risultati numerici per $k \in \{1000,2000,3000,4000,5000\}$ con \texttt{TIMES = 10}.}
\end{table}

\newpage

\section{Grafici}

\pgfplotsset{
    every axis/.append style={
        width=0.92\textwidth,
        height=0.42\textwidth,
        grid=both,
        xlabel={$k$},
        xtick={1000,2000,3000,4000,5000},
        scaled y ticks=true,
        legend style={at={(0.5,-0.22)}, anchor=north, legend columns=2},
        tick label style={font=\small},
        label style={font=\small},
        title style={font=\bfseries}
    }
}


\begin{figure}[H]
\centering
\begin{tikzpicture}
\begin{axis}[title={Avg fQ}, ylabel={Avg fQ}]
\addplot[blue, mark=*] coordinates {(1000,17350642347.4) (2000,18615897688.77) (3000,17038731132.6) (4000,20763157976.07) (5000,17864058550.1)};
\addlegendentry{QALS}
\addplot[red, dashed, mark=square*] coordinates {(1000,-714267512.7) (2000,-714267601.6) (3000,-714267631.8) (4000,-714267662) (5000,-714267662)};
\addlegendentry{NO QALS}
\end{axis}
\end{tikzpicture}
\caption{Andamento di Avg fQ al variare di $k$.}
\end{figure}


\begin{figure}[H]
\centering
\begin{tikzpicture}
\begin{axis}[title={fQ Min Found}, ylabel={fQ Min Found}]
\addplot[blue, mark=*] coordinates {(1000,8241260668.67) (2000,12244772130.67) (3000,12039123606.67) (4000,16020029113.67) (5000,6034515014.67)};
\addlegendentry{QALS}
\addplot[red, dashed, mark=square*] coordinates {(1000,-714267662) (2000,-714267662) (3000,-714267662) (4000,-714267662) (5000,-714267662)};
\addlegendentry{NO QALS}
\end{axis}
\end{tikzpicture}
\caption{Andamento di fQ Min Found al variare di $k$.}
\end{figure}


\begin{figure}[H]
\centering
\begin{tikzpicture}
\begin{axis}[title={Profit Found}, ylabel={Profit Found}]
\addplot[blue, mark=*] coordinates {(1000,3423) (2000,3481) (3000,4212) (4000,4249) (5000,2873)};
\addlegendentry{QALS}
\addplot[red, dashed, mark=square*] coordinates {(1000,2483) (2000,2483) (3000,2483) (4000,2483) (5000,2483)};
\addlegendentry{NO QALS}
\addplot[black, dotted, thick] coordinates {(1000,2679) (5000,2679)};
\addlegendentry{Profit ottimo}
\end{axis}
\end{tikzpicture}
\caption{Andamento del profitto trovato al variare di $k$.}
\end{figure}


\begin{figure}[H]
\centering
\begin{tikzpicture}
\begin{axis}[title={Weight Found}, ylabel={Weight Found}]
\addplot[blue, mark=*] coordinates {(1000,2275) (2000,2635) (3000,2618) (4000,2926) (5000,2041)};
\addlegendentry{QALS}
\addplot[red, dashed, mark=square*] coordinates {(1000,501) (2000,501) (3000,501) (4000,501) (5000,501)};
\addlegendentry{NO QALS}
\addplot[black, dotted, thick] coordinates {(1000,496) (5000,496)};
\addlegendentry{Weight ottimo}
\end{axis}
\end{tikzpicture}
\caption{Andamento del peso trovato al variare di $k$.}
\end{figure}


\begin{figure}[H]
\centering
\begin{tikzpicture}
\begin{axis}[title={Profit GAP}, ylabel={Profit GAP}]
\addplot[blue, mark=*] coordinates {(1000,-744) (2000,-802) (3000,-1533) (4000,-1570) (5000,-194)};
\addlegendentry{QALS}
\addplot[red, dashed, mark=square*] coordinates {(1000,196) (2000,196) (3000,196) (4000,196) (5000,196)};
\addlegendentry{NO QALS}
\addplot[black, dotted, thick] coordinates {(1000,0) (5000,0)};
\addlegendentry{GAP nullo}
\end{axis}
\end{tikzpicture}
\caption{Andamento del Profit GAP al variare di $k$.}
\end{figure}


\begin{figure}[H]
\centering
\begin{tikzpicture}
\begin{axis}[title={Weight GAP}, ylabel={Weight GAP}]
\addplot[blue, mark=*] coordinates {(1000,-1779) (2000,-2139) (3000,-2122) (4000,-2430) (5000,-1545)};
\addlegendentry{QALS}
\addplot[red, dashed, mark=square*] coordinates {(1000,-5) (2000,-5) (3000,-5) (4000,-5) (5000,-5)};
\addlegendentry{NO QALS}
\addplot[black, dotted, thick] coordinates {(1000,0) (5000,0)};
\addlegendentry{GAP nullo}
\end{axis}
\end{tikzpicture}
\caption{Andamento del Weight GAP al variare di $k$.}
\end{figure}


\begin{figure}[H]
\centering
\begin{tikzpicture}
\begin{axis}[title={n of fQ Min Found}, ylabel={n of fQ Min Found}]
\addplot[blue, mark=*] coordinates {(1000,1) (2000,1) (3000,1) (4000,1) (5000,1)};
\addlegendentry{QALS}
\addplot[red, dashed, mark=square*] coordinates {(1000,6) (2000,8) (3000,9) (4000,10) (5000,10)};
\addlegendentry{NO QALS}
\end{axis}
\end{tikzpicture}
\caption{Numero di volte in cui è stato trovato il minimo fQ al variare di $k$.}
\end{figure}


\newpage

\section{Soluzioni trovate}

In questa sezione sono riportati anche i vettori degli item selezionati per ogni valore di $k$.

\begin{table}[H]
\centering
\resizebox{\textwidth}{!}{%
\begin{tabular}{ccl}
\toprule
$k$ & Metodo & Items \\
\midrule
\textbf{1000} & QALS & \texttt{[1, 0, 1, 0, 0, 1, 0, 0, 0, 0, 0, 1, 1, 0, 0, 0, 1, 1, 0, 0, 0, 0, 0, 1, 1, 0, 0, 1, 1, 1]} \\
\textbf{1000} & NO QALS & \texttt{[0, 0, 0, 0, 0, 0, 0, 0, 0, 1, 0, 0, 0, 0, 0, 0, 1, 0, 0, 0, 0, 0, 0, 1, 1, 1, 0, 0, 1, 1]} \\
\textbf{2000} & QALS & \texttt{[1, 0, 0, 0, 0, 1, 1, 0, 1, 0, 1, 0, 0, 0, 0, 0, 1, 1, 0, 0, 1, 1, 1, 0, 1, 0, 0, 1, 1, 1]} \\
\textbf{2000} & NO QALS & \texttt{[0, 0, 0, 0, 0, 0, 0, 0, 0, 1, 0, 0, 0, 0, 0, 0, 1, 0, 0, 0, 0, 0, 0, 1, 1, 1, 0, 0, 1, 1]} \\
\textbf{3000} & QALS & \texttt{[0, 0, 0, 1, 1, 0, 0, 0, 0, 1, 1, 1, 0, 0, 1, 0, 0, 0, 0, 0, 1, 1, 1, 1, 1, 0, 0, 1, 1, 1]} \\
\textbf{3000} & NO QALS & \texttt{[0, 0, 0, 0, 0, 0, 0, 0, 0, 1, 0, 0, 0, 0, 0, 0, 1, 0, 0, 0, 0, 0, 0, 1, 1, 1, 0, 0, 1, 1]} \\
\textbf{4000} & QALS & \texttt{[0, 1, 0, 1, 0, 1, 1, 0, 1, 0, 0, 1, 0, 1, 0, 0, 1, 0, 0, 0, 1, 0, 1, 0, 0, 1, 1, 0, 1, 1]} \\
\textbf{4000} & NO QALS & \texttt{[0, 0, 0, 0, 0, 0, 0, 0, 0, 1, 0, 0, 0, 0, 0, 0, 1, 0, 0, 0, 0, 0, 0, 1, 1, 1, 0, 0, 1, 1]} \\
\textbf{5000} & QALS & \texttt{[0, 0, 0, 1, 0, 0, 0, 0, 0, 1, 0, 1, 0, 0, 0, 0, 0, 0, 1, 0, 1, 1, 1, 0, 1, 0, 1, 1, 1, 0]} \\
\textbf{5000} & NO QALS & \texttt{[0, 0, 0, 0, 0, 0, 0, 0, 0, 1, 0, 0, 0, 0, 0, 0, 1, 0, 0, 0, 0, 0, 0, 1, 1, 1, 0, 0, 1, 1]} \\
\bottomrule
\end{tabular}%
}
\caption{Vettori degli item trovati da QALS e NO QALS.}
\end{table}

\end{document}